\clearpage
\begin{explanationbox}
\textbf{\textcolor{CExplanation}{一、为什么会想到这个结论？}}
\begin{center}\begin{tikzpicture}[font=\small,>=stealth,line join=round]
\draw[gray!40,rounded corners=3pt] (0,0) rectangle (14.2,8.3);\node at (7.1,7.700000000000001) {例1：向右走2取倒数，向右走4才重复};
\begin{scope}[shift={(1.4956043956043956,2.2075757575757575)},x=0.6593406593406593cm,y=1.3939393939393938cm]\draw[->,gray] (-0.6,0)--(8.5,0) node[right]{$x$};\draw[->,gray] (0,-0.4)--(0,2.9) node[above]{$y$};\draw[blue,thick,solid,domain=-0.3:8,samples=180] plot(\x,{2^(sin(deg(pi*\x/2)))});\fill[orange!85!black] (1,2) circle(2.2pt);\fill[orange!85!black] (3,0.5) circle(2.2pt);\fill[orange!85!black] (5,2) circle(2.2pt);\fill[orange!85!black] (7,0.5) circle(2.2pt);\foreach \v in {1,3,5,7}{\node[below=5pt] at (\v,0){$\v$};}\end{scope}\node[anchor=west,align=left,text width=5.0cm] at (8.4,4.0) {$f=2^{\sin(\pi x/2)},\ a=2$\\[9pt]$f(1)=f(5)=2$\\[9pt]$f(3)=f(7)=\frac12$\\[9pt]$T=4$ 是一个周期};
\node[text width=12.6cm,align=center] at (7.1,.65) {按本题步长2和初值2绘制；曲线只作一条满足题设的示例。};
\end{tikzpicture}\end{center}



普通周期题会直接告诉你$f(x+4)=f(x)$，很容易看出“隔$4$就重复”。但有些题故意不让你直接看到重复，只告诉你
\[
 f(x+2)=\frac{1}{f(x)}.
\]
这时你看到的不是“值没变”，而是“值变成了倒数”。关键问题就变成：\textbf{再变一次会发生什么？}

\medskip

\end{explanationbox}
\clearpage
\begin{explanationbox}
\textbf{\textcolor{CExplanation}{二、先不看函数，玩一个初中生也懂的数字游戏}}
\begin{center}\begin{tikzpicture}[font=\small,>=stealth,line join=round]
\draw[gray!40,rounded corners=3pt] (0,0) rectangle (14.2,8.5);\node at (7.1,7.9) {把输出当作数：双曲线上的两点交换横纵坐标};
\node at (2.4,6.9) {$v=1/u$};\begin{scope}[shift={(2.4,4.4)},x=0.425cm,y=0.45cm]\draw[->,gray] (-4,0)--(4,0) node[right]{$u$};\draw[->,gray] (0,-4)--(0,4) node[above]{$v$};\draw[gray,thick,dashed,domain=-3.6:3.6,samples=160] plot(\x,{\x});\draw[blue,thick,solid,domain=-3.6:-0.28,samples=160] plot(\x,{1/\x});\draw[blue,thick,solid,domain=0.28:3.6,samples=160] plot(\x,{1/\x});\fill[orange!85!black] (2,0.5) circle(2pt);\fill[orange!85!black] (0.5,2) circle(2pt);\end{scope}\node[align=center,text width=4cm] at (2.4,1.7) {$2\to\frac12\to2$};\node at (7.050000000000001,6.9) {$v=-1/u$};\begin{scope}[shift={(7.050000000000001,4.4)},x=0.425cm,y=0.45cm]\draw[->,gray] (-4,0)--(4,0) node[right]{$u$};\draw[->,gray] (0,-4)--(0,4) node[above]{$v$};\draw[gray,thick,dashed,domain=-3.6:3.6,samples=160] plot(\x,{\x});\draw[violet,thick,solid,domain=-3.6:-0.28,samples=160] plot(\x,{-1/\x});\draw[violet,thick,solid,domain=0.28:3.6,samples=160] plot(\x,{-1/\x});\fill[orange!85!black] (2,-0.5) circle(2pt);\fill[orange!85!black] (-0.5,2) circle(2pt);\end{scope}\node[align=center,text width=4cm] at (7.050000000000001,1.7) {$2\to-\frac12\to2$};\node at (11.7,6.9) {$v=-6/u$};\begin{scope}[shift={(11.7,4.4)},x=0.425cm,y=0.45cm]\draw[->,gray] (-4,0)--(4,0) node[right]{$u$};\draw[->,gray] (0,-4)--(0,4) node[above]{$v$};\draw[gray,thick,dashed,domain=-3.6:3.6,samples=160] plot(\x,{\x});\draw[teal,thick,solid,domain=-3.6:-1.67,samples=160] plot(\x,{-6/\x});\draw[teal,thick,solid,domain=1.67:3.6,samples=160] plot(\x,{-6/\x});\fill[orange!85!black] (3,-2) circle(2pt);\fill[orange!85!black] (-2,3) circle(2pt);\end{scope}\node[align=center,text width=4cm] at (11.7,1.7) {$3\to-2\to3$};
\node[text width=12.6cm,align=center] at (7.1,.65) {横轴u为变换前的数，纵轴v为变换后的数；虚线v=u，两点关于它对称。};
\end{tikzpicture}\end{center}



从$2$出发，按“取倒数”的规则，每走一步就变一次：
\[
\boxed{2}\ \xrightarrow{\text{取倒数}}\
\boxed{\frac12}\ \xrightarrow{\text{再取倒数}}\
\boxed{2}\ \xrightarrow{\text{取倒数}}\
\boxed{\frac12}\ \cdots
\]
从$3$出发也一样：
\[
 3\ \longrightarrow\ \frac13\ \longrightarrow\ 3\ \longrightarrow\ \frac13\ \longrightarrow\cdots
\]
从$2$出发，若每次变成\textbf{负倒数}：
\[
\boxed{2}\ \longrightarrow\
\boxed{-\frac12}\ \longrightarrow\
\boxed{2}\ \longrightarrow\
\boxed{-\frac12}\ \cdots
\]
负号不会无限增加，因为“负倒数”做两次，前后两个负号正好消掉。

\medskip

\end{explanationbox}
\clearpage
\begin{explanationbox}
\textbf{\textcolor{CExplanation}{三、把数字游戏放回函数}}
\begin{center}\begin{tikzpicture}[font=\small,>=stealth,line join=round]
\draw[gray!40,rounded corners=3pt] (0,0) rectangle (14.2,8.3);\node at (7.1,7.700000000000001) {倒数型台阶：0、2处为2，1、3处为二分之一};
\begin{scope}[shift={(1.4296703296703297,2.323170731707317)},x=1.3186813186813187cm,y=1.1219512195121952cm]\draw[->,gray] (-0.25,0)--(4.3,0) node[right]{$x$};\draw[->,gray] (0,-0.6)--(0,3.5) node[above]{$y$};\draw[blue,thick,solid,domain=0:1,samples=180] plot(\x,{2});\draw[orange!85!black,thick,solid,domain=1:2,samples=180] plot(\x,{0.5});\draw[blue,thick,solid,domain=2:3,samples=180] plot(\x,{2});\draw[orange!85!black,thick,solid,domain=3:4,samples=180] plot(\x,{0.5});\fill[orange!85!black] (0,2) circle(2.2pt);\fill[orange!85!black] (1,0.5) circle(2.2pt);\fill[orange!85!black] (2,2) circle(2.2pt);\fill[orange!85!black] (3,0.5) circle(2.2pt);\draw[blue,fill=white] (1,2) circle(2.8pt);\draw[blue,fill=white] (2,0.5) circle(2.8pt);\draw[blue,fill=white] (3,2) circle(2.8pt);\draw[blue,fill=white] (4,0.5) circle(2.8pt);\foreach \v in {0,1,2,3,4}{\ifnum\v=0\node[below left=5pt] at (\v,-0.39999999999999997){$\v$};\else\node[below=5pt] at (\v,-0.39999999999999997){$\v$};\fi}\fill[orange!85!black] (4,2) circle(2.2pt);\end{scope}\node[anchor=west,align=left,text width=5.0cm] at (8.4,4.0) {$a=1,\ k=1$\\[9pt]$f(0)=f(2)=2$\\[9pt]$f(1)=f(3)=\frac12$\\[9pt]相邻值的乘积为1\\[9pt]整个图像每隔2重复};
\node[text width=12.6cm,align=center] at (7.1,.65) {实点取到，空圈不取；每段长度按本图标注的步长，隔两段重复。};
\end{tikzpicture}\end{center}



设$f(0)=2$，并且$f(x+1)=1/f(x)$，那么可立即填写：
\[
\begin{array}{c|ccccc}
 x &0&1&2&3&4\\ \hline
 f(x)&2&\frac12&2&\frac12&2
\end{array}
\]
每次输入$x$增加$1$，函数值就变一次；输入增加$2$，函数值就完整地走完“变过去、再变回来”。注意：\textbf{我们既不需要知道$f(0.3)$是多少，也不需要知道$f(x)$的具体公式}，就能确定$f(x+2)=f(x)$。

\medskip

\end{explanationbox}
\clearpage
\begin{explanationbox}
\textbf{\textcolor{CExplanation}{四、看图：走一步高度变了，走两步高度又一样}}
\begin{center}\begin{tikzpicture}[font=\small,>=stealth,line join=round]
\draw[gray!40,rounded corners=3pt] (0,0) rectangle (14.2,8.3);\node at (7.1,7.700000000000001) {倒数型台阶：0、2处为2，1、3处为二分之一};
\begin{scope}[shift={(1.4296703296703297,2.323170731707317)},x=1.3186813186813187cm,y=1.1219512195121952cm]\draw[->,gray] (-0.25,0)--(4.3,0) node[right]{$x$};\draw[->,gray] (0,-0.6)--(0,3.5) node[above]{$y$};\draw[blue,thick,solid,domain=0:1,samples=180] plot(\x,{2});\draw[orange!85!black,thick,solid,domain=1:2,samples=180] plot(\x,{0.5});\draw[blue,thick,solid,domain=2:3,samples=180] plot(\x,{2});\draw[orange!85!black,thick,solid,domain=3:4,samples=180] plot(\x,{0.5});\fill[orange!85!black] (0,2) circle(2.2pt);\fill[orange!85!black] (1,0.5) circle(2.2pt);\fill[orange!85!black] (2,2) circle(2.2pt);\fill[orange!85!black] (3,0.5) circle(2.2pt);\draw[blue,fill=white] (1,2) circle(2.8pt);\draw[blue,fill=white] (2,0.5) circle(2.8pt);\draw[blue,fill=white] (3,2) circle(2.8pt);\draw[blue,fill=white] (4,0.5) circle(2.8pt);\foreach \v in {0,1,2,3,4}{\ifnum\v=0\node[below left=5pt] at (\v,-0.39999999999999997){$\v$};\else\node[below=5pt] at (\v,-0.39999999999999997){$\v$};\fi}\fill[orange!85!black] (4,2) circle(2.2pt);\end{scope}\node[anchor=west,align=left,text width=5.0cm] at (8.4,4.0) {$a=1,\ k=1$\\[9pt]$f(0)=f(2)=2$\\[9pt]$f(1)=f(3)=\frac12$\\[9pt]相邻值的乘积为1\\[9pt]整个图像每隔2重复};
\node[text width=12.6cm,align=center] at (7.1,.65) {实点取到，空圈不取；每段长度按本图标注的步长，隔两段重复。};
\end{tikzpicture}\end{center}



下图给出两个真实的函数例子。让$a=1$，在每段长度为$1$的区间内取固定值：
\[
\begin{array}{c|cccc}
 \text{区间}&[0,1)&[1,2)&[2,3)&[3,4)\\ \hline
 \text{倒数型}&2&\frac12&2&\frac12\\
 \text{负倒数型}&2&-\frac12&2&-\frac12
\end{array}
\]

\end{explanationbox}\clearpage\begin{explanationbox}
\begin{center}\begin{tikzpicture}[font=\small,>=stealth,line join=round]
\draw[gray!40,rounded corners=3pt] (0,0) rectangle (14.2,8.3);\node at (7.1,7.700000000000001) {负倒数型台阶：跳过横轴，两次仍回到2};
\begin{scope}[shift={(1.4296703296703297,3.03)},x=1.3186813186813187cm,y=0.9199999999999999cm]\draw[->,gray] (-0.25,0)--(4.3,0) node[right]{$x$};\draw[->,gray] (0,-1.5)--(0,3.5) node[above]{$y$};\draw[blue,thick,solid,domain=0:1,samples=180] plot(\x,{2});\draw[orange!85!black,thick,solid,domain=1:2,samples=180] plot(\x,{-0.5});\draw[blue,thick,solid,domain=2:3,samples=180] plot(\x,{2});\draw[orange!85!black,thick,solid,domain=3:4,samples=180] plot(\x,{-0.5});\fill[orange!85!black] (0,2) circle(2.2pt);\fill[orange!85!black] (1,-0.5) circle(2.2pt);\fill[orange!85!black] (2,2) circle(2.2pt);\fill[orange!85!black] (3,-0.5) circle(2.2pt);\draw[blue,fill=white] (1,2) circle(2.8pt);\draw[blue,fill=white] (2,-0.5) circle(2.8pt);\draw[blue,fill=white] (3,2) circle(2.8pt);\draw[blue,fill=white] (4,-0.5) circle(2.8pt);\foreach \v in {0,1,2,3,4}{\ifnum\v=0\node[below left=5pt] at (\v,-1.3){$\v$};\else\node[below=5pt] at (\v,-1.3){$\v$};\fi}\fill[orange!85!black] (4,2) circle(2.2pt);\end{scope}\node[anchor=west,align=left,text width=5.0cm] at (8.4,4.0) {$a=1,\ k=-1$\\[9pt]$2\to-\frac12\to2$\\[9pt]相邻值的乘积为 $-1$\\[9pt]函数从来不取0\\[9pt]在整数处不连续};
\node[text width=12.6cm,align=center] at (7.1,.65) {实点取到，空圈不取；每段长度按本图标注的步长，隔两段重复。};
\end{tikzpicture}\end{center}

\textbf{读图时只看两件事：}蓝、橙线分别标出相邻两段高度；向右增加2后，整段重新出现。图中的“实心点”属于那一段，“空心点”不属于那一段，因此一个$x$仍然只对应一个函数值。

\medskip

\end{explanationbox}
\clearpage
\begin{explanationbox}
\textbf{\textcolor{CExplanation}{五、“平移”是输入在变，不是整条图像直接向右挪}}
\begin{center}\begin{tikzpicture}[font=\small,>=stealth,line join=round]
\draw[gray!40,rounded corners=3pt] (0,0) rectangle (14.2,8.3);\node at (7.1,7.700000000000001) {一次互为倒数，两次回到同一高度};
\begin{scope}[shift={(1.7923076923076926,2.0951612903225807)},x=1.153846153846154cm,y=1.4838709677419355cm]\draw[->,gray] (-0.6,0)--(4.6,0) node[right]{$x$};\draw[->,gray] (0,-0.3)--(0,2.8) node[above]{$y$};\draw[blue,thick,solid,domain=-0.3:4.3,samples=180] plot(\x,{2^(sin(deg(pi*\x)))});\fill[orange!85!black] (0.5,2) circle(2.2pt);\fill[orange!85!black] (1.5,0.5) circle(2.2pt);\fill[orange!85!black] (2.5,2) circle(2.2pt);\draw[gray,dashed] (-.5,1)--(4.4,1);\node[left] at (0,1){$1$};\foreach \v/\lab in {.5/\frac12,1.5/\frac32,2.5/\frac52}{\node[below=5pt] at (\v,0){$\lab$};}\end{scope}\node[anchor=west,align=left,text width=5.0cm] at (8.4,4.0) {$f=2^{\sin(\pi x)},\ a=1$\\[9pt]$f(\frac12)=2$\\[9pt]$f(\frac32)=\frac12$\\[9pt]$f(\frac52)=2$\\[9pt]正周期可取2};
\node[text width=12.6cm,align=center] at (7.1,.65) {三点高度是2、二分之一、2；倒数不同于把高度变为相反数。};
\end{tikzpicture}\end{center}



$f(x+1)=1/f(x)$的意思是：图上$x$位置的点是$(x,y)$，那么$x+1$位置的点应该是$\left(x+1,\frac1y\right)$。\textbf{只移动了横坐标还不够，纵坐标也要按规则改变。}第二次再走$1$，纵坐标才回到$y$。这才是“间隔$2$图像重复”的来源。

\medskip

\end{explanationbox}
\clearpage
\begin{explanationbox}
\textbf{\textcolor{CExplanation}{六、两个容易让人眼前一亮的联系}}
\begin{center}\begin{tikzpicture}[font=\small,>=stealth,line join=round]
\draw[gray!40,rounded corners=3pt] (0,0) rectangle (14.2,8.3);\node at (7.1,7.700000000000001) {证明三：面积相同、底边相同，所以高相同};
\draw[step=.5,gray!12] (3.4,2.1) grid (4.4,6.1);\draw[step=.5,gray!12] (9.2,2.1) grid (10.2,6.1);\filldraw[fill=blue!12,draw=blue,thick] (3.4,2.1) rectangle (4.4,6.1);\filldraw[fill=orange!12,draw=orange!85!black,thick] (9.2,2.1) rectangle (10.2,6.1);\node[above] at (3.9,6.2){面积1};\node[above] at (9.7,6.2){面积1};\node[left] at (3.3,4.1){$f(x)=2$};\node[right] at (10.3,4.1){$f(x+2)=2$};\node[below=6pt] at (3.9,2.1){$f(x+1)=\frac12$};\node[below=6pt] at (9.7,2.1){$f(x+1)=\frac12$};
\node[text width=12.6cm,align=center] at (7.1,.65) {图取 $a=k=1$；两矩形同尺度，共同底边 $f(x+1)\ne0$ 才能消去。};
\end{tikzpicture}\end{center}



\textbf{联系1：乘起来总是同一个数。}
条件$f(x+a)=k/f(x)$，换个写法就是
\[
\boxed{\ f(x)f(x+a)=k\ }.
\]
沿着数轴往右再走$a$，又有$f(x+a)f(x+2a)=k$。这两个乘积的\textbf{中间因子完全相同}，因此最左边和最右边的函数值必须相同。\textbf{有时“先把分式变成乘法”比反复套分式更快。}

\end{explanationbox}\clearpage\begin{explanationbox}
\begin{center}\begin{tikzpicture}[font=\small,>=stealth,line join=round]
\draw[gray!40,rounded corners=3pt] (0,0) rectangle (14.2,8.5);\node at (7.1,7.9) {同样满足加1取倒数，最小周期的情况却不同};
\node at (2.4,6.9) {$f\equiv1$};\begin{scope}[shift={(1.2230769230769232,3.0500000000000003)},x=0.8717948717948718cm,y=1.1250000000000002cm]\draw[->,gray] (-0.6,0)--(3.3,0) node[right]{$x$};\draw[->,gray] (0,-0.4)--(0,2.8) node[above]{$y$};\draw[blue,thick,solid,domain=-0.4:3,samples=160] plot(\x,{1});\end{scope}\node[align=center,text width=4cm] at (2.4,1.7) {任意正数都是周期\\没有最小正周期};\node at (7.050000000000001,6.9) {台阶：$2,\frac12$ 交替};\begin{scope}[shift={(5.8730769230769235,3.0500000000000003)},x=0.8717948717948718cm,y=1.1250000000000002cm]\draw[->,gray] (-0.6,0)--(3.3,0) node[right]{$x$};\draw[->,gray] (0,-0.4)--(0,2.8) node[above]{$y$};\draw[blue,thick,solid,domain=0:1,samples=160] plot(\x,{2});\draw[orange!85!black,thick,solid,domain=1:2,samples=160] plot(\x,{.5});\draw[blue,thick,solid,domain=2:3,samples=160] plot(\x,{2});\foreach \v/\h in {0/2,1/.5,2/2}{\fill[orange!85!black] (\v,\h) circle(2pt);}\foreach \v/\h in {1/2,2/.5,3/2}{\draw[blue,fill=white] (\v,\h) circle(2.5pt);}\fill[orange!85!black] (3,.5) circle(2pt);\end{scope}\node[align=center,text width=4cm] at (7.050000000000001,1.7) {整段每隔2重复\\最小正周期为2};\node at (11.7,6.9) {$g=2^{\sin(3\pi x)}$};\begin{scope}[shift={(10.523076923076923,3.0500000000000003)},x=0.8717948717948718cm,y=1.1250000000000002cm]\draw[->,gray] (-0.6,0)--(3.3,0) node[right]{$x$};\draw[->,gray] (0,-0.4)--(0,2.8) node[above]{$y$};\draw[violet,thick,solid,domain=-0.4:3,samples=160] plot(\x,{2^(sin(deg(3*pi*\x)))});\end{scope}\node[align=center,text width=4cm] at (11.7,1.7) {更早已经重复\\最小正周期为 $\frac23$};
\node[text width=12.6cm,align=center] at (7.1,.65) {三图同尺度，都有f(x+1)=1/f(x)；这条关系只能保证2是一个周期。};
\end{tikzpicture}\end{center}

\textbf{联系2：$2|a|$不一定是最小正周期。}
如果$f(x)\equiv1$，那么$f(x+a)=1/f(x)$对任意$a$都成立，但这个常数函数向右移动任何正长度都不变，所以不存在最小正周期。考试遇到“求最小正周期”，不能只根据这个关系就直接断定答案一定是$2|a|$。

\medskip

\end{explanationbox}
\clearpage
\begin{explanationbox}
\textbf{\textcolor{CExplanation}{七、原来倒数型也可以是一条平滑的曲线！}}
\begin{center}\begin{tikzpicture}[font=\small,>=stealth,line join=round]
\draw[gray!40,rounded corners=3pt] (0,0) rectangle (14.2,8.3);\node at (7.1,7.700000000000001) {一次互为倒数，两次回到同一高度};
\begin{scope}[shift={(1.7923076923076926,2.0951612903225807)},x=1.153846153846154cm,y=1.4838709677419355cm]\draw[->,gray] (-0.6,0)--(4.6,0) node[right]{$x$};\draw[->,gray] (0,-0.3)--(0,2.8) node[above]{$y$};\draw[blue,thick,solid,domain=-0.3:4.3,samples=180] plot(\x,{2^(sin(deg(pi*\x)))});\fill[orange!85!black] (0.5,2) circle(2.2pt);\fill[orange!85!black] (1.5,0.5) circle(2.2pt);\fill[orange!85!black] (2.5,2) circle(2.2pt);\draw[gray,dashed] (-.5,1)--(4.4,1);\node[left] at (0,1){$1$};\foreach \v/\lab in {.5/\frac12,1.5/\frac32,2.5/\frac52}{\node[below=5pt] at (\v,0){$\lab$};}\end{scope}\node[anchor=west,align=left,text width=5.0cm] at (8.4,4.0) {$f=2^{\sin(\pi x)},\ a=1$\\[9pt]$f(\frac12)=2$\\[9pt]$f(\frac32)=\frac12$\\[9pt]$f(\frac52)=2$\\[9pt]正周期可取2};
\node[text width=12.6cm,align=center] at (7.1,.65) {三点高度是2、二分之一、2；倒数不同于把高度变为相反数。};
\end{tikzpicture}\end{center}



别以为只有上面的“台阶图”才能满足倒数关系。把高中学过的正弦函数和指数函数连在一起，例如
\[
\boxed{\ f(x)=2^{\sin(\pi x)}\ },
\]
也有
\[
\begin{aligned}
 f(x+1)&=2^{\sin(\pi x+\pi)}
 =2^{-\sin(\pi x)}\\
 &=\frac1{2^{\sin(\pi x)}}=\frac1{f(x)}.
\end{aligned}
\]
因为$\sin(\pi x+\pi)=-\sin(\pi x)$，而\(2^{-u}=1/2^u\)。图像会像波浪一样平滑地起伏，却仍然\textbf{每走$2$就完整重复}。这说明周期来自\textbf{“走两次函数值回到原数”这个关系}，并不取决于图像是台阶还是波浪。


\medskip

\end{explanationbox}
\clearpage
\begin{explanationbox}
\textbf{\textcolor{CExplanation}{八、更深一层：负倒数型为什么很难画成一条不断开的曲线？}}
\begin{center}\begin{tikzpicture}[font=\small,>=stealth,line join=round]
\draw[gray!40,rounded corners=3pt] (0,0) rectangle (14.2,8.3);\node at (7.1,7.700000000000001) {负倒数型台阶：跳过横轴，两次仍回到2};
\begin{scope}[shift={(1.4296703296703297,3.03)},x=1.3186813186813187cm,y=0.9199999999999999cm]\draw[->,gray] (-0.25,0)--(4.3,0) node[right]{$x$};\draw[->,gray] (0,-1.5)--(0,3.5) node[above]{$y$};\draw[blue,thick,solid,domain=0:1,samples=180] plot(\x,{2});\draw[orange!85!black,thick,solid,domain=1:2,samples=180] plot(\x,{-0.5});\draw[blue,thick,solid,domain=2:3,samples=180] plot(\x,{2});\draw[orange!85!black,thick,solid,domain=3:4,samples=180] plot(\x,{-0.5});\fill[orange!85!black] (0,2) circle(2.2pt);\fill[orange!85!black] (1,-0.5) circle(2.2pt);\fill[orange!85!black] (2,2) circle(2.2pt);\fill[orange!85!black] (3,-0.5) circle(2.2pt);\draw[blue,fill=white] (1,2) circle(2.8pt);\draw[blue,fill=white] (2,-0.5) circle(2.8pt);\draw[blue,fill=white] (3,2) circle(2.8pt);\draw[blue,fill=white] (4,-0.5) circle(2.8pt);\foreach \v in {0,1,2,3,4}{\ifnum\v=0\node[below left=5pt] at (\v,-1.3){$\v$};\else\node[below=5pt] at (\v,-1.3){$\v$};\fi}\fill[orange!85!black] (4,2) circle(2.2pt);\end{scope}\node[anchor=west,align=left,text width=5.0cm] at (8.4,4.0) {$a=1,\ k=-1$\\[9pt]$2\to-\frac12\to2$\\[9pt]相邻值的乘积为 $-1$\\[9pt]函数从来不取0\\[9pt]在整数处不连续};
\node[text width=12.6cm,align=center] at (7.1,.65) {实点取到，空圈不取；每段长度按本图标注的步长，隔两段重复。};
\end{tikzpicture}\end{center}



负的$a$可以先改写成正步长的同一关系，因此只需讨论$a>0$。设$f(x+a)=-1/f(x)$。如果$f(x)$是正数，那么$f(x+a)$必为负数；反过来也一样。也就是说，隔$a$的两点一正一负。若曲线在这两点之间\textbf{不断开}，它就一定要经过横轴，即存在某处$f(x)=0$；但负倒数规则又要求函数值不能等于$0$，矛盾。因此：

\[
\boxed{\text{在整个实数范围内连续的实值函数，不可能处处满足}
\ f(x+a)=-\frac1{f(x)}.}
\]
这不是说负倒数型不存在，而是说\textbf{它若定义在全体实数上，就不可能处处连续}。上图中正值与负值之间的“跳变”正是一个很直观的例子。

\medskip

\end{explanationbox}
\clearpage
\begin{explanationbox}
\textbf{\textcolor{CExplanation}{九、考场最快识别法}}
\begin{center}\begin{tikzpicture}[font=\small,>=stealth,line join=round]
\draw[gray!40,rounded corners=3pt] (0,0) rectangle (14.2,8.3);\node at (7.1,7.700000000000001) {一次互为倒数，两次回到同一高度};
\begin{scope}[shift={(1.7923076923076926,2.0951612903225807)},x=1.153846153846154cm,y=1.4838709677419355cm]\draw[->,gray] (-0.6,0)--(4.6,0) node[right]{$x$};\draw[->,gray] (0,-0.3)--(0,2.8) node[above]{$y$};\draw[blue,thick,solid,domain=-0.3:4.3,samples=180] plot(\x,{2^(sin(deg(pi*\x)))});\fill[orange!85!black] (0.5,2) circle(2.2pt);\fill[orange!85!black] (1.5,0.5) circle(2.2pt);\fill[orange!85!black] (2.5,2) circle(2.2pt);\draw[gray,dashed] (-.5,1)--(4.4,1);\node[left] at (0,1){$1$};\foreach \v/\lab in {.5/\frac12,1.5/\frac32,2.5/\frac52}{\node[below=5pt] at (\v,0){$\lab$};}\end{scope}\node[anchor=west,align=left,text width=5.0cm] at (8.4,4.0) {$f=2^{\sin(\pi x)},\ a=1$\\[9pt]$f(\frac12)=2$\\[9pt]$f(\frac32)=\frac12$\\[9pt]$f(\frac52)=2$\\[9pt]正周期可取2};
\node[text width=12.6cm,align=center] at (7.1,.65) {三点高度是2、二分之一、2；倒数不同于把高度变为相反数。};
\end{tikzpicture}\end{center}



看到\(f(x+a)=\frac{k}{f(x)}\)（$k\ne0$），心里只想三个动作：\textbf{把$x$换成$x+a$；再把原式代入；最后找$f(x+2a)=f(x)$。}题目没要求写出函数表达式，就不要浪费时间去猜。
\end{explanationbox}
